
\documentclass{article}
\usepackage{colortbl}
\usepackage{makecell}
\usepackage{multirow}
\usepackage{supertabular}

\begin{document}

\newcounter{utterance}

\twocolumn

{ \footnotesize  \setcounter{utterance}{1}
\setlength{\tabcolsep}{0pt}
\begin{supertabular}{c@{$\;$}|p{.15\linewidth}@{}p{.15\linewidth}p{.15\linewidth}p{.15\linewidth}p{.15\linewidth}p{.15\linewidth}}

    \# & $\;$A & \multicolumn{4}{c}{Game Master} & $\;\:$B\\
    \hline 

    \theutterance \stepcounter{utterance}  

    & & \multicolumn{4}{p{0.6\linewidth}}{\cellcolor[rgb]{0.9,0.9,0.9}{%
	\makecell[{{p{\linewidth}}}]{% 
	  \tt {\tiny [A$\langle$GM]}  
	 Sie spielen ein kooperatives Verhandlungsspiel, bei dem Sie sich mit einem anderen Spieler darauf einigen müssen, wie eine Reihe von Gegenständen aufgeteilt werden soll.\\ \tt \\ \tt Die Regeln:\\ \tt (a) Sie und der andere Spieler erhalten eine Sammlung von Gegenständen. Jeder von Ihnen erhält außerdem eine geheime Wertfunktion, die angibt, wie viel Ihnen jede Art von Gegenstand wert ist.\\ \tt (b) Sie tauschen Nachrichten mit dem anderen Spieler aus, um zu vereinbaren, wer welche Gegenstände bekommt. Sie können jeweils maximal 5 Nachrichten senden oder das Spiel vorzeitig beenden, indem Sie jederzeit einen geheimen Vorschlag machen.\\ \tt (c) Jeder von euch wird aufgefordert, einen geheimen Vorschlag zu machen, in dem ihr die gewünschten Gegenstände in eckigen Klammern wie folgt angibt: "[Vorschlag: <Nummer> <Objektname>, <Nummer> <Objektname>, <...>]"\\ \tt (d) Wenn eure Vorschläge komplementär sind, d.h. es gibt genug Gegenstände, um beide Vorschläge zu erfüllen, erhält jeder Spieler eine Punktzahl, die sich aus der Summe der Werte für die Gegenstände ergibt, die er erhalten hat. Andernfalls erhalten Sie beide null Punkte.\\ \tt (e) Ihr Ziel ist es, die Summe aus Ihrer Punktzahl und der Punktzahl des anderen Spielers zu maximieren.\\ \tt \\ \tt Beginnen wir.\\ \tt \\ \tt Die Menge der verfügbaren Gegenstände ist:\\ \tt \\ \tt 1 Kartoffel, 1 Handschuh, 2 Lineale, 1 Pfeife, 2 Karten\\ \tt \\ \tt Deine geheime Wertfunktion ist:\\ \tt \\ \tt Kartoffel: 0, Handschuh: 1, Lineal: 3, Pfeife: 1, Karte: 1\\ \tt \\ \tt WICHTIG: Ihre Nachrichten werden, sofern es sich nicht um einen geheimen Vorschlag handelt, direkt an den anderen Spieler übermittelt, also fügen Sie keine Antwort auf die Regeln oder einen Text zur Ankündigung Ihrer Nachricht ein. Um einen geheimen Vorschlag zu machen, verwenden Sie das angegebene Format. Verwenden Sie keine eckigen Klammern, wenn Sie mit dem anderen Spieler kommunizieren, sonst wird dies als Ihr geheimer Vorschlag interpretiert.\\ \tt  
	  } 
	   } 
	   } 
	 & \\ 
 

    \theutterance \stepcounter{utterance}  

    & \multicolumn{4}{p{0.6\linewidth}}{\cellcolor[rgb]{0.8,1,0.9}{%
	\makecell[{{p{\linewidth}}}]{% 
	  \tt {\tiny [A$\rangle$GM]}  
	 Ich schlage vor, dass du die Kartoffel und die Pfeife behältst, und ich nehme beide Lineale und eine Karte. Der Handschuh und die zweite Karte können wir dann fair verteilen. Was meinst du? 
	  } 
	   } 
	   } 
	 & & \\ 
 

    \theutterance \stepcounter{utterance}  

    & & & \multicolumn{2}{p{0.3\linewidth}}{\cellcolor[rgb]{0.95,0.95,0.95}{%
	\makecell[{{p{\linewidth}}}]{% 
	  \tt {\tiny [GM$|$GM]}  
	 continue 
	  } 
	   } 
	   } 
	 & & \\ 
 

    \theutterance \stepcounter{utterance}  

    & & & \multicolumn{2}{p{0.3\linewidth}}{\cellcolor[rgb]{0.95,0.95,0.95}{%
	\makecell[{{p{\linewidth}}}]{% 
	  \tt {\tiny [GM$|$GM]}  
	 Ich schlage vor, dass du die Kartoffel und die Pfeife behältst, und ich nehme beide Lineale und eine Karte. Der Handschuh und die zweite Karte können wir dann fair verteilen. Was meinst du? 
	  } 
	   } 
	   } 
	 & & \\ 
 

    \theutterance \stepcounter{utterance}  

    & & \multicolumn{4}{p{0.6\linewidth}}{\cellcolor[rgb]{0.9,0.9,0.9}{%
	\makecell[{{p{\linewidth}}}]{% 
	  \tt {\tiny [GM$\rangle$B]}  
	 Sie spielen ein kooperatives Verhandlungsspiel, bei dem Sie sich mit einem anderen Spieler darauf einigen müssen, wie eine Reihe von Gegenständen aufgeteilt werden soll.\\ \tt \\ \tt Die Regeln:\\ \tt (a) Sie und der andere Spieler erhalten eine Sammlung von Gegenständen. Jeder von Ihnen erhält außerdem eine geheime Wertfunktion, die angibt, wie viel Ihnen jede Art von Gegenstand wert ist.\\ \tt (b) Sie tauschen Nachrichten mit dem anderen Spieler aus, um zu vereinbaren, wer welche Gegenstände bekommt. Sie können jeweils maximal 5 Nachrichten senden oder das Spiel vorzeitig beenden, indem Sie jederzeit einen geheimen Vorschlag machen.\\ \tt (c) Jeder von euch wird aufgefordert, einen geheimen Vorschlag zu machen, in dem ihr die gewünschten Gegenstände in eckigen Klammern wie folgt angibt: "[Vorschlag: <Nummer> <Objektname>, <Nummer> <Objektname>, <...>]"\\ \tt (d) Wenn eure Vorschläge komplementär sind, d.h. es gibt genug Gegenstände, um beide Vorschläge zu erfüllen, erhält jeder Spieler eine Punktzahl, die sich aus der Summe der Werte für die Gegenstände ergibt, die er erhalten hat. Andernfalls erhalten Sie beide null Punkte.\\ \tt (e) Ihr Ziel ist es, die Summe aus Ihrer Punktzahl und der Punktzahl des anderen Spielers zu maximieren.\\ \tt \\ \tt Beginnen wir.\\ \tt \\ \tt Die Menge der verfügbaren Gegenstände ist:\\ \tt \\ \tt 1 Kartoffel, 1 Handschuh, 2 Lineale, 1 Pfeife, 2 Karten\\ \tt \\ \tt Deine geheime Wertfunktion ist:\\ \tt \\ \tt Kartoffel: 7, Handschuh: 0, Lineal: 0, Pfeife: 3, Karte: 0\\ \tt \\ \tt WICHTIG: Ihre Nachrichten werden, sofern es sich nicht um einen geheimen Vorschlag handelt, direkt an den anderen Spieler übermittelt, also fügen Sie keine Antwort auf die Regeln oder einen Text zur Ankündigung Ihrer Nachricht ein. Um einen geheimen Vorschlag zu machen, verwenden Sie das angegebene Format. Verwenden Sie keine eckigen Klammern, wenn Sie mit dem anderen Spieler kommunizieren, sonst wird dies als Ihr geheimer Vorschlag interpretiert.\\ \tt \\ \tt \\ \tt Ich schlage vor, dass du die Kartoffel und die Pfeife behältst, und ich nehme beide Lineale und eine Karte. Der Handschuh und die zweite Karte können wir dann fair verteilen. Was meinst du? 
	  } 
	   } 
	   } 
	 & \\ 
 

    \theutterance \stepcounter{utterance}  

    & & & \multicolumn{4}{p{0.6\linewidth}}{\cellcolor[rgb]{1,0.85,0.72}{%
	\makecell[{{p{\linewidth}}}]{% 
	  \tt {\tiny [GM$\langle$B]}  
	 Ich könnte die Kartoffel und die Pfeife behalten, aber ich brauche auch mindestens eine Karte, weil ich sonst nur 10 Punkte hätte. Du bekommst die beiden Lineale und den Handschuh, ich nehme zusätzlich eine Karte — passt das? 
	  } 
	   } 
	   } 
	  \\ 
 

    \theutterance \stepcounter{utterance}  

    & & & \multicolumn{2}{p{0.3\linewidth}}{\cellcolor[rgb]{0.95,0.95,0.95}{%
	\makecell[{{p{\linewidth}}}]{% 
	  \tt {\tiny [GM$|$GM]}  
	 continue 
	  } 
	   } 
	   } 
	 & & \\ 
 

    \theutterance \stepcounter{utterance}  

    & & & \multicolumn{2}{p{0.3\linewidth}}{\cellcolor[rgb]{0.95,0.95,0.95}{%
	\makecell[{{p{\linewidth}}}]{% 
	  \tt {\tiny [GM$|$GM]}  
	 Ich könnte die Kartoffel und die Pfeife behalten, aber ich brauche auch mindestens eine Karte, weil ich sonst nur 10 Punkte hätte. Du bekommst die beiden Lineale und den Handschuh, ich nehme zusätzlich eine Karte — passt das? 
	  } 
	   } 
	   } 
	 & & \\ 
 

    \theutterance \stepcounter{utterance}  

    & & \multicolumn{4}{p{0.6\linewidth}}{\cellcolor[rgb]{0.9,0.9,0.9}{%
	\makecell[{{p{\linewidth}}}]{% 
	  \tt {\tiny [A$\langle$GM]}  
	 Ich könnte die Kartoffel und die Pfeife behalten, aber ich brauche auch mindestens eine Karte, weil ich sonst nur 10 Punkte hätte. Du bekommst die beiden Lineale und den Handschuh, ich nehme zusätzlich eine Karte — passt das? 
	  } 
	   } 
	   } 
	 & \\ 
 

    \theutterance \stepcounter{utterance}  

    & \multicolumn{4}{p{0.6\linewidth}}{\cellcolor[rgb]{0.8,1,0.9}{%
	\makecell[{{p{\linewidth}}}]{% 
	  \tt {\tiny [A$\rangle$GM]}  
	 Ich kann auf die Kartoffel verzichten, wenn du dafür die Pfeife behältst. Ich nehme beide Lineale, den Handschuh und eine Karte — das macht für mich 3+3+1+1=8 Punkte. Du bekommst die Kartoffel, die Pfeife und die zweite Karte — das sind 0+1+1=2 Punkte. Alternativ: Ich nehme nur ein Lineal, den Handschuh und beide Karten (3+1+1+1=6), du bekommst beide Lineale, die Kartoffel und die Pfeife (3+3+0+1=7). Das ergibt zusammen 13 statt 10. Was ist dir lieber? 
	  } 
	   } 
	   } 
	 & & \\ 
 

    \theutterance \stepcounter{utterance}  

    & & & \multicolumn{2}{p{0.3\linewidth}}{\cellcolor[rgb]{0.95,0.95,0.95}{%
	\makecell[{{p{\linewidth}}}]{% 
	  \tt {\tiny [GM$|$GM]}  
	 continue 
	  } 
	   } 
	   } 
	 & & \\ 
 

    \theutterance \stepcounter{utterance}  

    & & & \multicolumn{2}{p{0.3\linewidth}}{\cellcolor[rgb]{0.95,0.95,0.95}{%
	\makecell[{{p{\linewidth}}}]{% 
	  \tt {\tiny [GM$|$GM]}  
	 Ich kann auf die Kartoffel verzichten, wenn du dafür die Pfeife behältst. Ich nehme beide Lineale, den Handschuh und eine Karte — das macht für mich 3+3+1+1=8 Punkte. Du bekommst die Kartoffel, die Pfeife und die zweite Karte — das sind 0+1+1=2 Punkte. Alternativ: Ich nehme nur ein Lineal, den Handschuh und beide Karten (3+1+1+1=6), du bekommst beide Lineale, die Kartoffel und die Pfeife (3+3+0+1=7). Das ergibt zusammen 13 statt 10. Was ist dir lieber? 
	  } 
	   } 
	   } 
	 & & \\ 
 

    \theutterance \stepcounter{utterance}  

    & & \multicolumn{4}{p{0.6\linewidth}}{\cellcolor[rgb]{0.9,0.9,0.9}{%
	\makecell[{{p{\linewidth}}}]{% 
	  \tt {\tiny [GM$\rangle$B]}  
	 Ich kann auf die Kartoffel verzichten, wenn du dafür die Pfeife behältst. Ich nehme beide Lineale, den Handschuh und eine Karte — das macht für mich 3+3+1+1=8 Punkte. Du bekommst die Kartoffel, die Pfeife und die zweite Karte — das sind 0+1+1=2 Punkte. Alternativ: Ich nehme nur ein Lineal, den Handschuh und beide Karten (3+1+1+1=6), du bekommst beide Lineale, die Kartoffel und die Pfeife (3+3+0+1=7). Das ergibt zusammen 13 statt 10. Was ist dir lieber? 
	  } 
	   } 
	   } 
	 & \\ 
 

    \theutterance \stepcounter{utterance}  

    & & & \multicolumn{4}{p{0.6\linewidth}}{\cellcolor[rgb]{1,0.85,0.72}{%
	\makecell[{{p{\linewidth}}}]{% 
	  \tt {\tiny [GM$\langle$B]}  
	 Ich nehme die Kartoffel und die Pfeife, du bekommst beide Lineale, den Handschuh und beide Karten. Das ergibt für mich 7+3=10, für dich 0+0+0+0+0=0 — das geht nicht, du hast keine Punkte.  \\ \tt \\ \tt Alternativ: Ich nehme nur die Kartoffel, du bekommst die Pfeife, beide Lineale, den Handschuh und beide Karten. Dann habe ich 7, du hast 3 — zusammen 10.  \\ \tt Besser: Ich nehme die Kartoffel und eine Karte (7+0=7), du bekommst die Pfeife, beide Lineale, den Handschuh und die andere Karte (3+0+0+0+0=3). Zusammen 10.  \\ \tt \\ \tt Aber wenn du die Pfeife und die Kartoffel haben willst, und ich die restlichen vier Gegenstände — dann hast du 10, ich 0.  \\ \tt \\ \tt Mein Vorschlag: Du nimmst die Pfeife und eine Karte (3+0=3), ich nehme die Kartoffel und die andere Karte (7+0=7). Die beiden Lineale und den Handschuh schenken wir — der ist ja wertlos.  \\ \tt \\ \tt [ Vorschlag: 1 Kartoffel, 1 
	  } 
	   } 
	   } 
	  \\ 
 

    \theutterance \stepcounter{utterance}  

    & & & \multicolumn{2}{p{0.3\linewidth}}{\cellcolor[rgb]{0.95,0.95,0.95}{%
	\makecell[{{p{\linewidth}}}]{% 
	  \tt {\tiny [GM$|$GM]}  
	 continue 
	  } 
	   } 
	   } 
	 & & \\ 
 

    \theutterance \stepcounter{utterance}  

    & & & \multicolumn{2}{p{0.3\linewidth}}{\cellcolor[rgb]{0.95,0.95,0.95}{%
	\makecell[{{p{\linewidth}}}]{% 
	  \tt {\tiny [GM$|$GM]}  
	 Ich nehme die Kartoffel und die Pfeife, du bekommst beide Lineale, den Handschuh und beide Karten. Das ergibt für mich 7+3=10, für dich 0+0+0+0+0=0 — das geht nicht, du hast keine Punkte.  \\ \tt \\ \tt Alternativ: Ich nehme nur die Kartoffel, du bekommst die Pfeife, beide Lineale, den Handschuh und beide Karten. Dann habe ich 7, du hast 3 — zusammen 10.  \\ \tt Besser: Ich nehme die Kartoffel und eine Karte (7+0=7), du bekommst die Pfeife, beide Lineale, den Handschuh und die andere Karte (3+0+0+0+0=3). Zusammen 10.  \\ \tt \\ \tt Aber wenn du die Pfeife und die Kartoffel haben willst, und ich die restlichen vier Gegenstände — dann hast du 10, ich 0.  \\ \tt \\ \tt Mein Vorschlag: Du nimmst die Pfeife und eine Karte (3+0=3), ich nehme die Kartoffel und die andere Karte (7+0=7). Die beiden Lineale und den Handschuh schenken wir — der ist ja wertlos.  \\ \tt \\ \tt [ Vorschlag: 1 Kartoffel, 1 
	  } 
	   } 
	   } 
	 & & \\ 
 

    \theutterance \stepcounter{utterance}  

    & & \multicolumn{4}{p{0.6\linewidth}}{\cellcolor[rgb]{0.9,0.9,0.9}{%
	\makecell[{{p{\linewidth}}}]{% 
	  \tt {\tiny [A$\langle$GM]}  
	 Ich nehme die Kartoffel und die Pfeife, du bekommst beide Lineale, den Handschuh und beide Karten. Das ergibt für mich 7+3=10, für dich 0+0+0+0+0=0 — das geht nicht, du hast keine Punkte.  \\ \tt \\ \tt Alternativ: Ich nehme nur die Kartoffel, du bekommst die Pfeife, beide Lineale, den Handschuh und beide Karten. Dann habe ich 7, du hast 3 — zusammen 10.  \\ \tt Besser: Ich nehme die Kartoffel und eine Karte (7+0=7), du bekommst die Pfeife, beide Lineale, den Handschuh und die andere Karte (3+0+0+0+0=3). Zusammen 10.  \\ \tt \\ \tt Aber wenn du die Pfeife und die Kartoffel haben willst, und ich die restlichen vier Gegenstände — dann hast du 10, ich 0.  \\ \tt \\ \tt Mein Vorschlag: Du nimmst die Pfeife und eine Karte (3+0=3), ich nehme die Kartoffel und die andere Karte (7+0=7). Die beiden Lineale und den Handschuh schenken wir — der ist ja wertlos.  \\ \tt \\ \tt [ Vorschlag: 1 Kartoffel, 1 
	  } 
	   } 
	   } 
	 & \\ 
 

    \theutterance \stepcounter{utterance}  

    & \multicolumn{4}{p{0.6\linewidth}}{\cellcolor[rgb]{0.8,1,0.9}{%
	\makecell[{{p{\linewidth}}}]{% 
	  \tt {\tiny [A$\rangle$GM]}  
	 [ Vorschlag: 1 Pfeife, 1 Karte ] 
	  } 
	   } 
	   } 
	 & & \\ 
 

    \theutterance \stepcounter{utterance}  

    & & & \multicolumn{2}{p{0.3\linewidth}}{\cellcolor[rgb]{0.95,0.95,0.95}{%
	\makecell[{{p{\linewidth}}}]{% 
	  \tt {\tiny [GM$|$GM]}  
	 proposal 
	  } 
	   } 
	   } 
	 & & \\ 
 

    \theutterance \stepcounter{utterance}  

    & & & \multicolumn{2}{p{0.3\linewidth}}{\cellcolor[rgb]{0.95,0.95,0.95}{%
	\makecell[{{p{\linewidth}}}]{% 
	  \tt {\tiny [GM$|$GM]}  
	 [0, 0, 0, 1, 1] 
	  } 
	   } 
	   } 
	 & & \\ 
 

    \theutterance \stepcounter{utterance}  

    & & \multicolumn{4}{p{0.6\linewidth}}{\cellcolor[rgb]{0.9,0.9,0.9}{%
	\makecell[{{p{\linewidth}}}]{% 
	  \tt {\tiny [GM$\rangle$B]}  
	 Der andere Spieler hat einen geheimen Vorschlag gemacht. Jetzt bist du an der Reihe, einen eigenen Vorschlag zu machen.\\ \tt  
	  } 
	   } 
	   } 
	 & \\ 
 

    \theutterance \stepcounter{utterance}  

    & & & \multicolumn{4}{p{0.6\linewidth}}{\cellcolor[rgb]{1,0.85,0.72}{%
	\makecell[{{p{\linewidth}}}]{% 
	  \tt {\tiny [GM$\langle$B]}  
	 [ Vorschlag: 1 Kartoffel, 1 Pfeife ] 
	  } 
	   } 
	   } 
	  \\ 
 

    \theutterance \stepcounter{utterance}  

    & & & \multicolumn{2}{p{0.3\linewidth}}{\cellcolor[rgb]{0.95,0.95,0.95}{%
	\makecell[{{p{\linewidth}}}]{% 
	  \tt {\tiny [GM$|$GM]}  
	 proposal 
	  } 
	   } 
	   } 
	 & & \\ 
 

    \theutterance \stepcounter{utterance}  

    & & & \multicolumn{2}{p{0.3\linewidth}}{\cellcolor[rgb]{0.95,0.95,0.95}{%
	\makecell[{{p{\linewidth}}}]{% 
	  \tt {\tiny [GM$|$GM]}  
	 [1, 0, 0, 1, 0] 
	  } 
	   } 
	   } 
	 & & \\ 
 

    \theutterance \stepcounter{utterance}  

    & & & \multicolumn{2}{p{0.3\linewidth}}{\cellcolor[rgb]{0.95,0.95,0.95}{%
	\makecell[{{p{\linewidth}}}]{% 
	  \tt {\tiny [GM$|$GM]}  
	 [[0, 0, 0, 1, 1], [1, 0, 0, 1, 0]] 
	  } 
	   } 
	   } 
	 & & \\ 
 

\end{supertabular}
}

\end{document}
